\subsection{函数演算的连续性}

在本节中，我们考虑函数演算关于两个变量的连续性质。

对星代数 A 及 A 的正规元素 \(a\) ，以及对函数 \(f\) ——在 \(\mathbf{C}\) 的包含 \(\mathrm{Sp}_{\mathrm{A}}(a)\) 的子集上连续——我们用 \(f(a)\) 表示把 \(a\) 的连续函数演算应用于 \(f\) 在 \(a\) 的谱上的限制所得到的元素。

命题 10. --- 设 A 是含幺星代数。设 U 是 \(\mathbf{C}\) 中的相对紧开集。用 \(\Omega_{\mathrm{n}}\) 表示 A 中满足 \(\mathrm{Sp}_{\mathrm{A}}(a)\) 含于 U 的正规元素所成的集。映射 \((f,a)\mapsto f(a)\) ——从 \(\mathcal{C}(\overline{\mathrm{U}})\times \Omega_{\mathrm{n}}\) 到 A 中——是连续的。

用 \(q\) 表示映射 \((f, a) \mapsto f(a)\) ——从 \(\mathcal{C}(\overline{\mathrm{U}}) \times \Omega_{\mathrm{n}}\) 到 A 中。函数演算映射 \(f \mapsto f(a)\) （\(a \in \Omega_{\mathrm{n}}\)）所成的族在 \(\mathcal{L}(\mathcal{C}(\overline{\mathrm{U}}); \mathrm{A})\) 中是等度连续的，因为其中每个都是范数 \(\leqslant 1\) 的连续线性映射。因此，为证明断言，只须验证：对每个 \(f \in \mathcal{C}(\overline{\mathrm{U}})\) ，从 \(\Omega_{\mathrm{n}}\) 到 A 中的、由 \(a \mapsto f(a)\) 所定义的映射是连续的（TG, X, p. 13，推论 3）。

设 \(\mathcal{A}\) 是由 \(f\in\mathcal{C}(\overline{\mathrm{U}})\) （使映射 \(a\mapsto f(a)\) ——从 \(\Omega_{\mathrm{n}}\) 到 A 中——连续）组成的集；要证明的是 \(\mathcal{A}=\mathcal{C}(\overline{\mathrm{U}})\) 。

集 \(\mathcal{A}\) 是 \(\mathcal{C}(\overline{\mathrm{U}})\) 的含幺对合子代数。它包含 \(\overline{\mathrm{U}}\) 上的恒等函数，因而分离点。于是，它在 \(\mathcal{C}(\overline{\mathrm{U}})\) 中稠密（TG, X, p. 39，命题 7）。我们来证明它是闭的。

设 \((f_{\mathrm{n}})\) 是 \(\mathcal{A}\) 中收敛于 \(f \in \mathcal{C}(\overline{\mathrm{U}})\) 的序列。设 \(\varepsilon > 0\) ，并选取 \(n \in \mathbf{N}\) 使得 \(\|f - f_{n}\|_{\infty} \leqslant \varepsilon/4\) 。对每个 \((a_{1}, a_{2}) \in \Omega_{\mathrm{n}}^{2}\) ，我们有

\[
\begin{array}{c} \| f (a _ {1}) - f (a _ {2}) \| \leqslant 2 \| f - f _ {n} \| _ {\infty} + \| f _ {n} (a _ {1}) - f _ {n} (a _ {2}) \| \\ \leqslant \frac {\varepsilon}{2} + \| f _ {n} (a _ {1}) - f _ {n} (a _ {2}) \|. \end{array}
\]

由于 \(f_n\) 属于 \(\mathcal{A}\) ，存在 \(a_1\) 在 \(\Omega_n\) 中的邻域 V，使得 \(\| f_n(a_1) - f_n(a_2) \| \leqslant \varepsilon / 2\) 对所有 \(a_2 \in V\) 成立。于是，我们有 \(\| f(a_1) - f(a_2) \| \leqslant \varepsilon\) （对所有 \(a_2 \in V\) ）。因而映射 \(a \mapsto f(a)\) 在 \(a_1\) 处连续；断言得证。

推论 1. --- 设 A 是含幺星代数，\(A_{n}\) 是 A 的正规元素所成的集。给空间 \(\mathcal{C}(\mathbf{C})\) 赋以紧收敛拓扑。映射 \((f,a)\mapsto f(a)\) ——从 \(\mathcal{C}(\mathbf{C})\times\mathrm{A}_{\mathrm{n}}\) 到 A 中——是连续的。

设 \((f_{0}, a_{0}) \in \mathcal{C}(\mathbf{C}) \times \mathrm{A}_{\mathrm{n}}\) 。设 U 是 \(a_{0}\) 的谱的相对紧邻域。设 V 是 A 中满足 \(\mathrm{Sp}_{\mathrm{A}}(a) \subset \mathrm{U}\) 的正规元素 a 所成的集；它是 \(a_{0}\) 在 \(A_{n}\) 中的开邻域（I, p. 76，命题 10）。对每个 \(a \in V\) 及每个函数 \(f \in \mathcal{C}(\mathbf{C})\) ，我们有 \(f(a) = (f|\overline{\mathrm{U}})(a)\) 。由于映射 \(f \mapsto \|f|\overline{U}\|_{\infty}\) 是赋以紧收敛拓扑的空间 \(\mathcal{C}(\mathbf{C})\) 上的连续半范数，映射 \((f, a) \mapsto f(a)\) 在 \((f_{0}, a_{0})\) 处的连续性由命题 10 得出。

推论 2. --- 设 E 是复希尔伯特空间，\(\mathcal{L}(\mathrm{E})_{\mathrm{n}}\) 是 E 的正规自同态所成的集。给空间 \(\mathcal{C}(\mathbf{C})\) 赋以紧收敛拓扑。从 \(\mathcal{C}(\mathbf{C}) \times \mathcal{L}(\mathrm{E})_{\mathrm{n}}\) 到 \(\mathcal{L}(\mathrm{E})\) 中的、由 \((f, u) \mapsto f(u)\) 所定义的映射是连续的。

